% !PRINTED 82
\weberpage{82}
%<<CONT
\gesperrt{dann und nur dann besteht, wenn zugleich die
Congruenzen}
\weq{5}{x_1\equiv 0\ (\mathrm{mod}\ i_1),\quad x_2\equiv 0\ (\mathrm{mod}\ i_2),\quad \dots\ x_\nu\equiv 0\ (\mathrm{mod}\ i_\nu)}
\gesperrt{erfüllt sind.}

Hiernach ist also $J$ die kleinste positive Zahl, die für alle $b$ der Congruenz
\[
b^J\equiv 1\ (\mathrm{mod}\ m)
\]
genügt, oder die kleinste positive Zahl, für die für alle Systeme der $\beta$
\weq{6}{J\beta_1\equiv 0\ (\mathrm{mod}\ c_1),\quad J\beta_2\equiv 0\ (\mathrm{mod}\ c_2),\quad \dots\ J\beta_\mu\equiv 0\ (\mathrm{mod}\ c_\mu).}

Daraus ergeben sich die ersten Schlüsse über die Zusammensetzung der Zahl $m$.

a) \gesperrt{Eine ungerade Primzahl $q$, die nicht in $J$ enthalten ist, kann nur einfach in $m$ aufgehen.}

Denn ist $m$ durch $q^\varkappa$ theilbar, so ist eine der Grössen $c$, etwa $c_1=\varphi(q^\varkappa)=q^{\varkappa-1}(q-1)$, also wenn $\varkappa>1$ ist, so ist $c_1$ noch durch $q$ theilbar, und aus (6) folgt dann, dass alle $\beta_1$ durch $q$ theilbar sein müssten, was der Voraussetzung IV. widerspricht.

b) \gesperrt{Eine ungerade Primzahl $q$, die in $J$ aufgeht, kann höchstens einmal mehr in $m$, als in $J$ enthalten sein.}

Denn man kann ebenso schliessen, dass, wenn $c_1=q^{\varkappa-1}(q-1)$, und $J$ nicht durch $q^{\varkappa-1}$ theilbar ist, alle $\beta_1$ durch $q$ theilbar sein müssten.

c) \gesperrt{Ist $J$ ungerade, so muss auch $m$ ungerade sein.}

Denn ist $m$ durch eine Potenz von 2, also mindestens durch 4 theilbar, so ist der der Zahl 2 entsprechende Indexmodul gerade, und die entsprechenden $\beta$ müssten nach (6) alle gerade sein, entgegen der Forderung IV.

d) \gesperrt{Ist $J$ durch eine Potenz von 2 theilbar, so kann $m$ den Factor 2 höchstens zweimal öfter enthalten, als $J$.}

Denn ist $m$ durch $2^\lambda$ theilbar und $\lambda>2$, so sind zwei der Indexmoduln
\[
c_1=2,\quad c_2=2^{\lambda-2}.
\]
%CONT>>
